% Compiles with any TeX distribution: pdflatex/xelatex/lualatex abstract.tex
% Uses only base packages (geometry, hyperref) — portable, one page.
\documentclass[11pt]{article}

\usepackage[a4paper,margin=1in]{geometry}
\usepackage[hidelinks]{hyperref}
\setlength{\parskip}{0.5em}
\setlength{\parindent}{0pt}

\title{\vspace{-1.5em}Reharvester: A Local-First System for Interactive
Research-Literature Discovery}
\author{Yongyuth Chuankhuntod and Anirach Mingkhwan}
\date{}

\begin{document}
\maketitle
\vspace{-2em}

\begin{abstract}
\noindent
Keeping pace with a fast-growing research literature is increasingly difficult, and
the cloud services that promise to help are opaque, non-portable, and demand that
sensitive reading lists leave the user's machine. We present \emph{Reharvester}, a
local-first system that turns a keyword or abstract query into an interactive map of a
research field without any external dependency. Reharvester harvests papers, assembles
them into a navigable knowledge graph, surfaces emerging keyword trends, and exposes the
underlying corpus through a wiki-linked reader backed by retrieval-augmented generation
over the graph. The entire pipeline runs on the user's hardware and degrades gracefully
when compute or models are unavailable, making literature exploration reproducible,
private, and air-gappable.
\end{abstract}

\section{Introduction}
The volume of published research outpaces any individual's capacity to read, and
existing discovery tools centralize both the data and the reasoning in the cloud. This
is a poor fit for privacy-sensitive or offline settings, and it hides the ranking and
retrieval logic from the researcher. Reharvester reframes literature discovery as a
local, inspectable pipeline: every stage---ingestion, graph construction, trend
analysis, and question answering---executes on the user's own machine and remains open
to inspection.

\section{Approach}
Reharvester couples four stages. \textbf{Ingestion} gates a harvest on a meaningful
query---at least five keywords or a short abstract---then crawls matching papers while
streaming progress and a live terminal log to the user. \textbf{Graph construction}
renders the resulting corpus as a WebGL knowledge graph in which node size reflects
PageRank centrality, five interchangeable layouts reveal different structure, one-hop
highlighting traces neighborhoods, and a research-gap overlay marks sparsely connected
regions worth exploring. \textbf{Trend analysis} ranks keywords by the regression slope
of their frequency over a user-adjustable multi-year window, distinguishing durable
themes from noise. \textbf{Corpus access} pairs a searchable, open-access-filterable
paper list with an Obsidian-style wiki pane whose \texttt{[[wikilinks]]} select and
center the corresponding graph node; a local language model answers questions through
retrieval-augmented generation grounded in the graph.

\section{Implementation and Results}
The system is implemented as a Go harvesting server and a Next.js~16 single-page
interface using Reagraph for WebGL rendering and Zustand for client-global state that
survives navigation, so crawls and selections persist across views. A built-in mock mode
ships a seeded 450-node corpus, a simulated crawl stream, and trend and wiki fixtures,
letting the full interface run with no backend---useful for demonstration and offline
evaluation. The renderer scales to a ten-thousand-node stress mode, automatically
shedding labels and animation as node counts grow, and resilience is built in throughout:
health polling with automatic recovery, a WebGL fallback path, and an explicit
model-unreachable indicator when the local model is down. Together these choices deliver
a private, portable, and responsive tool for exploring a research field end to end.

\vspace{1em}
\noindent\textbf{Keywords:} literature discovery; knowledge graph; local-first;
retrieval-augmented generation; research-gap analysis; keyword trends.

\end{document}
